\def\level{Première}  % Classe à droite
\def\course{NSI}
\def\chaptername{Conditions IF} % Titre au centre
\def\docname{AP03-S}        % Identifiant à gauche

\pagestyle{cours_FP}
\makeheader{} % Affiche le bandeau


\begin{easybox}{Les structures conditionnelles IF}

\begin{minipage}{0.32\linewidth}
	\begin{tikzpicture}[node distance=1.5cm]
		% Nœuds
		\node (instruction1) [process] at (0, 0){Instruction 1};
		\node (decision) [decision] at (0, -2){Condition};
		\node (instruction2) [process] at (0, -4.5){Instruction 2};
		\node (instruction3) [process] at (0, -6) {Instruction 3};
	
		\coordinate (midpoint1) at ($(instruction2.south)!0.5!(instruction3.north)$);
		\coordinate (point1) at (2, -4.5);
		% Flèches
		\draw [arrow] (instruction1) -- (decision);
		\draw [arrow] (decision.south) -- node[anchor=west] {True} (instruction2.north);
		\draw [arrow] (instruction2) -- (instruction3);
		\draw [arrow] (decision.east) -| node[anchor=east, xshift=1.2cm, yshift=-2cm] {False}(point1);
		\draw [arrow] (point1) |- (midpoint1);
	\end{tikzpicture}	
	\end{minipage} \hfill
	\begin{minipage}{0.65\linewidth}
		\begin{minipage}{0.55\linewidth}
	\begin{scratch}
    \blocklook{mettre \ovalvariable{total} à \ovalvariable{10}}
\blockif{si \booloperator{\ovalmove{total} > \ovalnum{0}} alors}
{\blocklook{mettre \ovalvariable{total} à \ovaloperator{\ovalvariable{total} + \ovalnum{1}}}
}
\blocklook{Afficher \ovalvariable{total}}
\end{scratch}		
		\end{minipage}\hfill
		\begin{minipage}{0.43\linewidth}
	\begin{lstlisting}[numbers=none]
# En python
Instruction_1
if condition:
	Instruction_2
Instruction_3
	\end{lstlisting}
    
    \begin{lstlisting}[numbers=none]
# En python
total = 10
if total > 0:
	total = total + 1
print(total)
	\end{lstlisting}
		\end{minipage}
	
	
	\begin{itemize}
		\item l'\verb+Instruction 1+ est effectuée
		\item SI la \verb+Condition+ est \verb+True+ alors l'\verb+Instruction 2+ est effectuée
		\item SINON on va directement à l'\verb+Instruction 3+
		\item Dans tous les cas l'\verb+Instruction 3+ est effectuée
	\end{itemize}
	
	La \verb+Condition+ est un test qui est évalué en python par un booléen.
	
	\end{minipage}

\end{easybox}


	\begin{easybox}{Les structures conditionnelles IF -- ELSE}{}
		\begin{minipage}{0.4\linewidth}
			\begin{tikzpicture}[node distance=1.5cm]
				% Nœuds
				\node (instruction1) [process] at (0, 0){instruction 1};
				\node (decision) [decision] at (0, -2){Condition};
				\node (instruction2) [process] at (3.5, -4){instruction 3};
				\node (instruction3) [process] at (0, -4) {instruction 2};
				\node (instruction4) [process] at (0, -6) {instruction 4};
				
				\coordinate (midpoint1) at ($(instruction3.south)!0.5!(instruction4.north)$);
					% Flèches
				\draw [arrow] (instruction1) -- (decision);
				\draw [arrow] (decision.east) -| node[anchor=south] {False} (instruction2.north);
				\draw [arrow] (decision.south) -- node[anchor=east, xshift=-0.5cm, yshift=0.3cm] {True} (instruction3.north);
				\draw [arrow] (instruction3) -- (instruction4);
				\draw [arrow] (instruction2) |- (midpoint1);
		
			\end{tikzpicture}		
		\end{minipage}\hfill
		\begin{minipage}{0.58\linewidth}
			\begin{minipage}{0.48\linewidth}
	\begin{scratch}
\blockifelse{si \booloperator{\ovalmove{x} >= \ovalnum{0}} alors}
{\blocklook{mettre \ovalvariable{x} à {\ovalvariable{x ** 2} }}}
{\blocklook{mettre \ovalvariable{x} à {\ovalvariable{x ** 3} }}}
\end{scratch}	
			\end{minipage}\hfill
			\begin{minipage}{0.48\linewidth}
	\begin{lstlisting}[numbers=none]
# En python
	
instruction_1
	
if condition:
	instruction_2
else:
	instruction_3
instruction_4
	\end{lstlisting}		
			\end{minipage}
				
				
	\begin{itemize}
		\item l'\verb+instruction 1+ est effectuée
		\item SI la \verb+Condition+ est \verb+True+ alors l'\verb+instruction 2+ est effectuée
		\item SINON l'\verb+instruction 3+ est effectuée
		\item Dans tous les cas l'\verb+instruction 4+ est effectuée
	\end{itemize}
		\end{minipage}
			
			
	\end{easybox}
	
	\medskip

		\begin{exemple}[Les structures conditionnelles]{}
			
        \begin{minipage}{0.48\linewidth}
            \begin{CodePythontexAlt}[Largeur=1\linewidth,Centre,Lignes=false,PremLigne=1]{}
def plus_grand2(a, b):
    ''' Renvoie le plus grand des 
    deux nombres a et b '''
	if a > b:
		return a
	else:
		return b

maximum = plus_grand2(5, 10)
print(maximum) # Affiche 10
	\end{CodePythontexAlt}
        \end{minipage}\hfill{}
        \begin{minipage}{0.48\linewidth}
            \begin{CodePythontexAlt}[Largeur=1\linewidth,Centre,Lignes=false,PremLigne=1]{}
def plus_grand3(a, b, c):
    ''' Renvoie le plus grand des 
    trois nombres a, b et c '''
    if a > b and a > c:
		return a
	elif b > a and b > c:
		return b
	else:
		return c

maximum = plus_grand3(5, 10, 7)
print(maximum) # Affiche 10
		\end{CodePythontexAlt}
        \end{minipage}
        
	
	
	
		
	
	
	\end{exemple}	
